% !TeX program = xelatex
% 实变函数（Real Analysis）笔记（LaTeX 版）
% 依据手写笔记整理：共六章（集合论动机、Lebesgue 测度、可测函数、Lebesgue 积分、微分与不定积分、Lp 空间）。
\documentclass[11pt]{ctexart}
\usepackage[a4paper,margin=2.3cm]{geometry}
\usepackage{amsmath,amssymb,amsthm,mathtools}
\usepackage{booktabs}
\usepackage{array}
\usepackage[colorlinks=true,linkcolor=blue!55!black,urlcolor=blue!55!black]{hyperref}

\setcounter{secnumdepth}{-2}
\setcounter{tocdepth}{1}

\theoremstyle{definition}
\newtheorem*{definition}{定义}
\newtheorem*{proposition}{命题}
\newtheorem*{theorem}{定理}
\newtheorem*{lemma}{引理}
\newtheorem*{corollary}{推论}
\newtheorem*{example}{例}
\newtheorem*{remark}{注}
\newtheorem*{exercise}{习题}
\theoremstyle{remark}
\renewcommand{\proofname}{证明}

\newcommand{\R}{\mathbb{R}}
\newcommand{\Z}{\mathbb{Z}}
\newcommand{\N}{\mathbb{N}}

% 章标题（手工编号，与小节编号 N.x 对应）
\newcommand{\chap}[1]{%
  \clearpage
  \vspace*{0.5em}
  {\centering\LARGE\bfseries #1\par}
  \vspace{0.4em}
  \addcontentsline{toc}{section}{#1}}
\newcommand{\secn}[1]{\subsection*{#1}\addcontentsline{toc}{section}{#1}\markboth{#1}{}}

\title{\heiti 实变函数\\[0.3em]\large 课堂笔记（LaTeX 整理版）}
\author{}
\date{2026 年 10 月}

\begin{document}
\maketitle
\vspace{-3em}

\begin{abstract}
\noindent
本笔记整理自《实变函数》课程手写笔记，共六章：第一章实分析研究动机与集合论，第二章 Lebesgue 测度，
第三章可测函数及其性质，第四章 Lebesgue 积分，第五章微分与不定积分，第六章 $L^p$ 空间。
记号约定：$m$ 表示 Lebesgue 测度，$m^*$ 表示外测度，$\mathcal{M}$ 为 Lebesgue 可测集全体，
$\mathcal{B}(\mathbb{R}^n)$ 为 Borel $\sigma$-代数；$\chi_E$ 为示性函数；a.e. 表示几乎处处；
$R[a,b]$ 表示 $[a,b]$ 上 Riemann 可积函数全体；BV、AC 分别表示有界变差与绝对连续函数类；
$C_c$ 表示紧支撑连续函数。电子版（Markdown 一节一页）见个人主页课程笔记。
\end{abstract}

\newpage
\tableofcontents
\newpage

\chap{第一章　实分析研究动机与集合论}

本章从 Riemann 积分的三大缺陷出发，说明引入新积分理论的必要性；随后建立集合的对等与基数理论（可数集、不可数集、连续统），并引入集合列的极限、开集与闭集、Borel 集与 Cantor 集等预备知识，为第二章 Lebesgue 测度做准备。

\secn{1.1 Riemann 积分的局限性}

\begin{proposition}
$f \in R[a,b] \iff f$ 在 $[a,b]$ 上几乎处处连续（Lebesgue 判据）。
\end{proposition}

回忆积分学基本定理（FTC）：$F(x) = \int_a^x f(t)\,dt$，$F' = f$。Riemann 积分理论存在三方面的问题：

（1）某些简单函数并非 Riemann 可积。

\begin{example}[Dirichlet 函数]
$D(x) = \chi_{\mathbb{Q} \cap [0,1]}(x) \notin R[0,1]$：任意分割下振幅 $\omega \ge 1$。
\end{example}

（2）Newton–Leibniz 公式有额外限制：需要 $f'$ 可积或连续；存在 Volterra 型反例（函数可微、$f'$ 有界，但 $f'$ 非 R-可积）。

（3）积分与极限的交换一般不成立。

\begin{example}
$f_n = n \cdot \chi_{(0,1/n)}$：$\int_0^1 f_n = 1$，但 $f_n \to 0$ a.e. 于 $[0,1]$，$\int_0^1 0\,dm = 0$。
\end{example}

\begin{quote}
由此引出本课程的主线：\textbf{扩大可积函数类 + 建立更好的极限定理}。Newton–Leibniz 公式成立的"最小条件"将在第五章（绝对连续函数）中彻底解决。
\end{quote}

\secn{1.2 测度思想：分割值域}

\begin{definition}[测度的公理化愿望]
把区间长度推广为集合的"测度"：$m((a,b)) = b-a$，并希望 $m$ 满足非负性、可数可加性、平移不变性。
\end{definition}

Lebesgue 的关键想法：\textbf{分割值域而非定义域}。

\begin{itemize}

\item R-积分：分割定义域 $[a,b] = \bigcup_k I_k$，作黎曼和；

\item L-积分：分割值域 $f(E) \subseteq \bigcup_j [y_j, y_{j+1})$，取水平集 $E_j = f^{-1}([y_j, y_{j+1}))$，

\end{itemize}

\[
\int_E f \, dm = \sup_{\text{分割}} \sum_j \xi_j \, m(E_j), \qquad \xi_j \in [y_j, y_{j+1}).
\]

对比：\textbf{L-积分 = 分值域 + 可测集测度 + 上确界；R-积分 = 分定义域 + 区间长度 + 黎曼和}。这要求水平集 $E_j$ "可测"、$m$ 具备可数可加性——这正是第二章要建立的理论。

\secn{1.3 对等与基数}

\begin{definition}[对等]
$A \sim B \iff$ 存在双射 $f: A \to B$。对等是等价关系。
\end{definition}

\begin{example}
$\mathbb{N} \sim 2\mathbb{N}$（$n \mapsto 2n$）；$\mathbb{N} \sim \mathbb{Z}$；$(0,1) \sim \mathbb{R}$（如 $x \mapsto \tan\frac{\pi}{2}(2x-1)$）。
\end{example}

\begin{theorem}[1.1（Cantor–Bernstein）]
若 $A \sim B_1 \subset B$ 且 $B \sim A_1 \subset A$，则 $A \sim B$。
\end{theorem}

\begin{proof}
记 $f: A \to B_1$、$g: B \to A_1$ 为双射。令 $A_0 = A \setminus g(B)$，递归定义 $A_{n+1} = g(f(A_n))$，并令 $C = \bigcup_{n\ge 0} A_n$。构造

\[
h(x) = \begin{cases} f(x), & x \in C \\ g^{-1}(x), & x \notin C \end{cases}
\]

验证 $h$ 为 $A \to B$ 的双射。
\end{proof}

\secn{1.4 可数集与不可数集}

\begin{definition}[可数集]
$A \sim \mathbb{N}$；有限集与可数集统称至多可数集。
\end{definition}

\begin{theorem}[1.2]
$\mathbb{Q}$ 是可数集。
\end{theorem}

\begin{proof}
把 $\mathbb{Q}$ 写成按 $\vert p\vert + q$ 分层的既约分数之并，逐层枚举；或用 $(m,n) \mapsto 2^m 3^n$ 的单射嵌入 $\mathbb{N}$。
\end{proof}

\begin{theorem}[1.3]
可数个至多可数集之并仍至多可数。
\end{theorem}

\begin{proof}
列表法 + 对角线数法。
\end{proof}

\begin{corollary}
以下集合均可数：$\mathbb{Z}$、$\mathbb{Q}^n$、代数数全体。
\end{corollary}

\begin{theorem}[1.4（Cantor）]
$\mathbb{R}$ 是不可数集。
\end{theorem}

\begin{proof}
对角线法：设 $[0,1)$ 可数，列举为 $x_1, x_2, \ldots$（十进制展开），构造 $y$ 使其第 $n$ 位小数与 $x_n$ 的第 $n$ 位不同，则 $y$ 不在列表中，矛盾。
\end{proof}

\begin{corollary}
超越数存在且不可数；$[0,1] \sim (0,1) \sim \mathbb{R}$。
\end{corollary}

\secn{1.5 基数运算与连续统}

\begin{theorem}[1.5]
$\aleph_0 + \aleph_0 = \aleph_0$；对任意有限 $n$，$n \cdot \aleph_0 = \aleph_0$；$\aleph_0^2 = \aleph_0$。
\end{theorem}

\begin{theorem}[1.6（无最大基数）]
对任意集合 $A$，$\vert \mathcal{P}(A) \vert > \vert A \vert$；特别地 $2^{\aleph_0} > \aleph_0$。
\end{theorem}

\begin{proof}
显然 $\vert A \vert \le \vert \mathcal{P}(A) \vert$（单射 $a \mapsto \{a\}$）。若存在双射 $f: A \to \mathcal{P}(A)$，取 $B = \{x \in A : x \notin f(x)\}$，则 $B$ 无原像——对角线法。
\end{proof}

\begin{theorem}[1.7]
$\mathbb{R} \sim 2^{\mathbb{N}}$（连续统的刻画）。
\end{theorem}

\begin{proof}
二进制展开给出 $2^{\mathbb{N}} \to [0,1]$ 的满射；对展开式的二义性（可数集）作修正后用 Cantor–Bernstein。
\end{proof}

\begin{corollary}
$\mathbb{R}^n \sim \mathbb{R}$（$n \ge 1$）。
\end{corollary}

\begin{remark}
连续统假设 CH（"没有基数严格介于 $\aleph_0$ 与 $2^{\aleph_0}$ 之间"）独立于 ZFC（Gödel 1938 / Cohen 1963）。
\end{remark}

\begin{example}
$(0,1) \sim [0,1)$：取

\[
f(x) = \begin{cases} \frac{1}{n-1}, & x = \frac{1}{n},\ n \ge 2 \\ x, & \text{其他} \end{cases}
\]
\end{example}

（平移一个可数点集）。

\begin{theorem}[1.8]
无穷集必与自己的某个真子集对等；反之，与某真子集对等的集合是无穷集。
\end{theorem}

\secn{1.6 实数集的序结构与上确界}

\begin{theorem}[1.9（确界原理）]
非空有上界的数集必有上确界。
\end{theorem}

\begin{proof}
由 Dedekind 分割构造，或由单调有界数列极限存在（区间套）。
\end{proof}

\begin{remark}[Archimedes 性质]
$\forall \varepsilon > 0$，$\exists n \in \mathbb{N}$ 使 $\frac{1}{n} < \varepsilon$；有理数在 $\mathbb{R}$ 中稠密。
\end{remark}

\secn{1.7 集合列的极限}

\begin{definition}[上、下极限集]
\[
\varlimsup_{n \to \infty} E_n = \bigcap_{n=1}^{\infty} \bigcup_{k \ge n} E_k, \qquad \varliminf_{n \to \infty} E_n = \bigcup_{n=1}^{\infty} \bigcap_{k \ge n} E_k.
\]
\end{definition}

若二者相等，记其公共值为 $\lim_n E_n$。

\begin{example}
$E_n = [0,\, 1 + \frac{(-1)^n}{n}]$：$\varlimsup E_n = [0,1]$，$\varliminf E_n = [0,1)$。
\end{example}

\begin{definition}[示性函数]
$\chi_E(x) = 1$（$x \in E$），$0$（$x \notin E$）。
\end{definition}

\begin{proposition}
$\chi_{\varlimsup E_n} = \varlimsup \chi_{E_n}$，$\chi_{\varliminf E_n} = \varliminf \chi_{E_n}$（集合极限 = 函数逐点极限）。
\end{proposition}

\begin{proposition}[测度连续性的伏笔]
若 $E_n \nearrow E$（递增），则 $\chi_{E_n} \to \chi_E$；长度函数相应地满足 $m(E) = \lim_n m(E_n)$——严格的表述见第二章（测度的下连续性）。
\end{proposition}

\secn{1.8 开集与闭集}

\begin{definition}[开集]
$E \subseteq \mathbb{R}^n$ 为开集，若 $\forall x \in E$，$\exists \delta > 0$ 使开球 $B(x,\delta) \subseteq E$。
\end{definition}

\begin{definition}[闭集]
$E$ 为闭集，若 $E^c$ 为开集；等价地，$E$ 包含其全部极限点（$E = \overline{E}$）。
\end{definition}

\begin{theorem}[1.10]
任意多个开集之并为开集；有限多个开集之交为开集。对偶地：任意多个闭集之交为闭集；有限多个闭集之并为闭集。
\end{theorem}

\begin{theorem}[1.11（Borel 集）]
开集全体生成的 $\sigma$-代数 = 闭集全体生成的 $\sigma$-代数，记为 $\mathcal{B}(\mathbb{R}^n)$，其元素称为 Borel 集。
\end{theorem}

\begin{example}
$G_\delta$ 集（可数个开集之交）、$F_\sigma$ 集（可数个闭集之并）都是 Borel 集。
\end{example}

\secn{1.9 直线上开集的构造与 Cantor 集}

\begin{theorem}[1.12]
直线上的开集 $G$ 可写为至多可数个互不相交的开区间之并：$G = \bigcup_k I_k$（$I_k$ 称为构成区间）。
\end{theorem}

\begin{proof}
对 $x \in G$，令 $I_x$ 为包含 $x$ 且含于 $G$ 的极大开区间（所有此类区间之并）。不同的 $I_x$ 或相等或不交，且每个这样的区间含一个有理数，故至多可数个。
\end{proof}

\begin{remark}
$\mathbb{R}^n$（$n \ge 2$）中的开集没有如此简单的区间分解。
\end{remark}

\begin{definition}[Cantor 三分集]
从 $[0,1]$ 出发，逐次删除每个区间的中间三分之一开区间，余集记为 $C$。
\end{definition}

\begin{theorem}[1.13（Cantor 集的性质）]
$C$ 是闭的、无处稠密的完全集（没有孤立点的闭集），且 $C \sim \{0,1\}^{\mathbb{N}}$（基数连续统），但其 Lebesgue 测度为 $0$：删除区间总长

\[
\sum_{k=1}^{\infty} 2^{k-1} \cdot 3^{-k} = 1.
\]
\end{theorem}

\begin{proof}[证明思路]
三进制展开：$x \in C \iff x$ 有只含数字 $0,2$ 的三进制展开，由此 $C \sim 2^{\mathbb{N}}$；完全性由 $C = \bigcap_k C_k$（有限个闭区间之并）与端点结构得出。
\end{proof}

\begin{quote}
\textbf{本章小结}：Riemann 积分的三大缺陷（可积函数类太小、N–L 公式条件过强、极限定理缺失）$\Rightarrow$ 需要新的测度与积分理论；工具准备：对等与基数（Cantor–Bernstein、可数/不可数）、集合列极限、开集结构与 Borel 集、Cantor 集（重要的反例源泉）。
\end{quote}

\chap{第二章　Lebesgue 测度}

本章给 $\mathbb{R}^n$ 中的集合赋予"长度/体积"：从外测度出发，用 Carathéodory 条件选出可测集类，证明它是 $\sigma$-代数且测度在其上可数可加；随后讨论可测集的结构定理、Vitali 不可测集、一般测度的扩张（Carathéodory 定理与 Lebesgue–Stieltjes 测度）、完备化与乘积测度预备。

\secn{2.1 外测度}

\begin{definition}[方体]
$I = \prod_{k=1}^n (a_k, b_k)$，体积 $\mathrm{vol}(I) = \prod_{k=1}^n (b_k - a_k)$。
\end{definition}

\begin{definition}[Lebesgue 外测度]
对 $E \subseteq \mathbb{R}^n$，

\[
m^*(E) = \inf\left\{ \sum_{j=1}^{\infty} \mathrm{vol}(I_j) : E \subseteq \bigcup_{j=1}^{\infty} I_j,\ I_j \text{ 为开方体} \right\}
\]
\end{definition}

约定 $m^*(\emptyset) = 0$，取值于 $[0, +\infty]$。

\begin{proposition}[2.1（外测度的基本性质）]

\end{proposition}

（i）非负性：$m^*(E) \ge 0$，$m^*(\emptyset) = 0$；

（ii）单调性：$E \subseteq F \Rightarrow m^*(E) \le m^*(F)$；

（iii）次可数可加性：$m^*\left(\bigcup_j E_j\right) \le \sum_j m^*(E_j)$；

（iv）平移不变性：$m^*(E + x) = m^*(E)$。

\begin{proof}
（iii）：任给 $\varepsilon > 0$，对每个 $E_j$ 取开方体覆盖使 $\sum_k \mathrm{vol}(I_j^{(k)}) \le m^*(E_j) + \varepsilon/2^j$；合并所有方体即得 $E$ 的覆盖，令 $\varepsilon \to 0$。
\end{proof}

\begin{example}
可数集的外测度为 $0$：$m^*(\mathbb{Q} \cap [0,1]) = 0$（单点覆盖长度 $\varepsilon/2^j$）。
\end{example}

\begin{example}[区间的外测度 = 长度]
$m^*((a,b)) = b - a$。
\end{example}

\begin{proof}
"$\le$"：用自身覆盖。"$\ge$"：对 $[a+\delta,\, b-\delta]$ 的任意开覆盖取有限子覆盖（紧性），令 $\delta \to 0$。
\end{proof}

\begin{proposition}[2.2（距离外测度性质）]
若 $d(E, F) > 0$，则 $m^*(E \cup F) = m^*(E) + m^*(F)$。
\end{proposition}

\begin{proof}
把 $E \cup F$ 的覆盖分解为直径小于 $d(E,F)$ 的方体，则每个方体至多与 $E$、$F$ 之一相交，分成两族分别覆盖。
\end{proof}

\begin{quote}
外测度在全体集合上\textbf{不是}可数可加的（见 2.3 的 Vitali 集），因此需要把"可测"的集合从一般集合中甄别出来。
\end{quote}

\secn{2.2 可测集与可数可加性}

\begin{definition}[Carathéodory 条件]
$E \subseteq \mathbb{R}^n$ 称为 Lebesgue 可测，若对任意 $T \subseteq \mathbb{R}^n$，

\[
m^*(T) = m^*(T \cap E) + m^*(T \setminus E).
\]

直观：$E$ 可测 $\iff$ $E$ 能把任一集合的外测度"精确分割"成两部分。可测集全体记为 $\mathcal{M}$。
\end{definition}

\begin{example}
$\emptyset$、$\mathbb{R}^n$、外测度为零的集合均可测。
\end{example}

\begin{theorem}[2.1]
$\mathcal{M}$ 是 $\sigma$-代数（从而包含所有 Borel 集），且限制 $m := m^*\vert_{\mathcal{M}}$ 在 $\mathcal{M}$ 上可数可加，构成完备测度。
\end{theorem}

\begin{proof}[证明思路]
补集封闭：在 Carathéodory 条件中交换 $E$ 与 $E^c$ 的角色。有限并封闭：对 $T$ 用条件拆分两次。可数并封闭：对 $T \cap \bigcup_j E_j$ 归纳 + $\varepsilon/2^j$ 技巧，再由次可加性取极限；"$\ge$"方向先对有限并取等号，再用递增集列的逼近。
\end{proof}

\begin{corollary}
开集、闭集、$G_\delta$ 集、$F_\sigma$ 集均可测。
\end{corollary}

\begin{theorem}[2.2（可数可加性）]
若 $\{E_j\}$ 是两两不交的可测集，则

\[
m\left(\bigcup_{j=1}^{\infty} E_j\right) = \sum_{j=1}^{\infty} m(E_j).
\]
\end{theorem}

\begin{proof}
"$\le$"为次可加性；"$\ge$"先对有限并证明（对 $\bigcup_{j\le N} E_j$ 用 Carathéodory 条件，$T = \bigcup_j E_j$），再令 $N \to \infty$。
\end{proof}

\begin{corollary}[测度的连续性]
若 $E_j \nearrow E$，则 $m(E) = \lim_j m(E_j)$；若 $E_j \searrow E$ 且 $m(E_1) < \infty$，则 $m(E) = \lim_j m(E_j)$。
\end{corollary}

\secn{2.3 不可测集的存在性}

\begin{theorem}[2.3（Vitali）]
存在 $\mathbb{R}$ 的子集不是 Lebesgue 可测的。
\end{theorem}

\begin{proof}
在 $[0,1)$ 上定义等价关系 $x \sim y \iff x - y \in \mathbb{Q}$，由选择公理从每个等价类中取一点组成代表集 $V$。

设 $\{r_j\} = \mathbb{Q} \cap [-1, 1)$，则平移族 $\{V + r_j\}$ 两两不交，且

\[
[0,1) \subseteq \bigcup_j (V + r_j) \subseteq [-1, 2).
\]

若 $V$ 可测，由平移不变性与可数可加性：

\[
1 \le m\left(\textstyle\bigcup_j (V + r_j)\right) = \sum_j m(V) \le 3,
\]

但 $\sum_j m(V)$ 当 $m(V) = 0$ 时为 $0$、当 $m(V) > 0$ 时为 $+\infty$，两种情形均矛盾。
\end{proof}

\begin{remark}
证明用到选择公理；"所有集合都可测"与选择公理不相容（Solovay 模型）。
\end{remark}

\secn{2.4 可测集的结构}

\begin{theorem}[2.4（结构定理）]
对 $E \subseteq \mathbb{R}^n$，以下等价：
\end{theorem}

（1）$E$ 可测；

（2）$\forall \varepsilon > 0$，存在开集 $G \supseteq E$ 使 $m^*(G \setminus E) < \varepsilon$；

（3）存在 $G_\delta$ 集 $G \supseteq E$ 使 $m(G \setminus E) = 0$；

（4）存在 $F_\sigma$ 集 $F \subseteq E$ 使 $m(E \setminus F) = 0$。

\begin{proof}[证明思路]
（1）$\Rightarrow$（2）：先把 $E$ 截成有界块（$E_k = E \cap B_k$），对每块用外测度定义取方体覆盖作开集，误差 $\varepsilon/2^k$ 求和。（2）$\Rightarrow$（3）：取 $G_k$ 使 $m^*(G_k \setminus E) < 1/k$，令 $G = \bigcap_k G_k$。（3）$\Rightarrow$（1）：$E = G \setminus (G \setminus E)$，零测集可测。（4）：对 $E^c$ 用（3）。
\end{proof}

\begin{corollary}[可测集的分解]
$E$ 可测 $\iff$ $E = B \cup N$，其中 $B$ 为 Borel 集、$m(N) = 0$（或 $E = B \setminus N$）。
\end{corollary}

\begin{quote}
直观：\textbf{可测集 = Borel 集整体 + 零测集的修饰}。
\end{quote}

\secn{2.5 平移不变性与正则性}

\begin{proposition}
若 $E$ 可测，则 $m(E + x) = m(E)$（$x \in \mathbb{R}^n$）。
\end{proposition}

\begin{theorem}[2.5（正则性）]
$E$ 可测时，

\[
m(E) = \inf\{m(G) : G \text{开} \supseteq E\} = \sup\{m(K) : K \text{紧} \subseteq E\}.
\]
\end{theorem}

\begin{proof}[证明思路]
外测度定义直接给外正则（开覆盖）；内正则：$E \cap \overline{B}_k \nearrow E$，对有界情形用结构定理（4）中闭集逼近 + 紧性。
\end{proof}

\begin{example}
Cantor 集 $C$ 的测度：$C$ 是零测集（第一章已算删除总长为 1，或 $C \subseteq G_k$，$m(G_k) = (2/3)^k \to 0$）。
\end{example}

\secn{2.6 Steiner 对称化}

（几何应用：等周不等式等的测度论预备。）

\begin{definition}[关于超平面的对称化]
设 $E \subseteq \mathbb{R}^n$，$H$ 为超平面。把每条与 $H$ 垂直的弦"中点移到 $H$ 上"所得的集合记为 $E^S$。
\end{definition}

\begin{proposition}
$m(E^S) = m(E)$（可由 Fubini 定理逐方向验证，或作为 Brunn–Minkowski 不等式的推论）。
\end{proposition}

\begin{theorem}[2.6（等周不等式的预备）]
若 $E$ 有界可测，则存在对称化序列使 $m(E_k^S) = m(E)$ 恒成立，而边界"周长"不增；极限图形是球——圆解决等周问题。
\end{theorem}

\secn{2.7 一般测度与 Carathéodory 扩张定理}

\begin{definition}[代数上的预测度]
设 $\mathcal{A}$ 是 $X$ 上的代数（对有限并、补封闭），$\mu_0 : \mathcal{A} \to [0, \infty]$ 满足 $\mu_0(\emptyset) = 0$ 且在 $\mathcal{A}$ 内可数可加，称为预测度。
\end{definition}

\begin{theorem}[2.7（Carathéodory 扩张定理）]
$\mu_0$ 可以扩张为 $\sigma(\mathcal{A})$ 上的测度 $\mu$；若 $\mu_0$ 在 $\mathcal{A}$ 上 $\sigma$-有限，则扩张唯一。
\end{theorem}

\begin{proof}[证明思路]
令 $\mu^*(E) = \inf\{\sum_j \mu_0(A_j) : E \subseteq \bigcup_j A_j,\ A_j \in \mathcal{A}\}$，验证 $\mu^*$ 是外测度且每个 $A \in \mathcal{A}$ 满足 Carathéodory 条件；唯一性用单调类/$\lambda$-$\pi$ 系论证。
\end{proof}

\begin{example}

\end{example}

\begin{itemize}

\item Lebesgue 测度 = "长度"函数在方体（半开矩形）代数上的扩张；

\item 乘积测度（2.10）= 矩形上 $\mu \times \nu$ 的扩张。

\end{itemize}

\secn{2.8 Lebesgue–Stieltjes 测度}

\begin{definition}
设 $F: \mathbb{R} \to \mathbb{R}$ 单调递增、右连续。在形如 $(a, b]$ 的区间上定义 $\mu_F((a,b]) = F(b) - F(a)$，再作 Carathéodory 扩张。
\end{definition}

\begin{proposition}
$\mu_F$ 是 Borel 测度，且 $\mu_F((-\infty, x]) = F(x)$。
\end{proposition}

\begin{example}[跳跃函数]
$F$ 为阶梯函数时 $\mu_F$ 是纯点测度：$\mu_F(\{a\}) = F(a) - F(a^-)$（原子质量 = 跳跃高度）。
\end{example}

\begin{theorem}[2.8]
Borel 测度与单调递增右连续函数（相差常数意义下）一一对应。
\end{theorem}

\begin{quote}
概率论的联系：分布函数 $F$ $\leftrightarrow$ 概率测度 $\mu_F$（$\mu_F(\mathbb{R}) = 1$）。
\end{quote}

\secn{2.9 完备化}

\begin{definition}[完备性]
测度空间 $(X, \mathcal{F}, \mu)$ 完备：$\mu(N) = 0$、$N' \subseteq N$ $\Rightarrow$ $N' \in \mathcal{F}$。
\end{definition}

\begin{proposition}
Lebesgue 测度空间 $(\mathbb{R}^n, \mathcal{M}, m)$ 是完备的；但 Borel $\sigma$-代数上的 $m$ 不完备（存在含于零测 Borel 集的非 Borel 子集，见 2.11）。
\end{proposition}

\begin{theorem}[2.9]
任何 $\sigma$-有限测度空间可以唯一完备化：$\overline{\mathcal{F}} = \{B \cup N : B \in \mathcal{F},\ N \subseteq \text{零测集}\}$。
\end{theorem}

\begin{remark}
完备化不改变可积函数的积分值（零测集上的修改无影响）；$\mathcal{M} = \overline{\mathcal{B}(\mathbb{R}^n)}$。
\end{remark}

\secn{2.10 乘积测度预备}

（为第四章 Fubini 定理做准备。）

\begin{definition}[可测矩形]
$A \times B$（$A \in \mathcal{X}$，$B \in \mathcal{Y}$），定义 $(\mu \times \nu)(A \times B) = \mu(A)\,\nu(B)$，再由 Carathéodory 扩张得到乘积测度。
\end{definition}

\begin{proposition}[截口的可测性]
设 $E \in \mathcal{X} \otimes \mathcal{Y}$，则每个截口 $E_y = \{x : (x, y) \in E\} \in \mathcal{X}$；且 $\mu \times \nu$ $\sigma$-有限时，$y \mapsto \mu(E_y)$ 可测。
\end{proposition}

\begin{example}
$\mathbb{R}^2$ 上的 Lebesgue 测度 = $m \times m$；$\sigma$-有限性对截口函数性质不可省。
\end{example}

\secn{2.11 可测集类的大小与 Cantor 函数}

\begin{proposition}
$\mathcal{M}$（Lebesgue 可测集全体）的基数为 $2^{2^{\aleph_0}}$，而 Borel 集全体的基数只有 $2^{\aleph_0}$：\textbf{存在 Lebesgue 可测的非 Borel 集}（且"大多数"可测集不是 Borel 集）。
\end{proposition}

\begin{definition}[Cantor 函数 / 魔鬼楼梯]
在三进制删除区间上依次取常数 $1/2, 1/4, \ldots$，在 Cantor 集 $C$ 上连续延拓，所得 $G: [0,1] \to [0,1]$ 连续、递增、满射。
\end{definition}

\textbf{性质.} $G' = 0$ a.e. 但 $G$ 非常数（与 Newton–Leibniz 公式的关系见第五章）。

\textbf{应用（非 Borel 可测集的构造）.} 令 $H(x) = G(x) + x$，则 $H$ 严格递增、连续，把零测的 Cantor 集映为正测度集 $H(C)$（$m(H(C)) = 1$）。取 $H(C)$ 的不可测子集 $W$，则 $H^{-1}(W) \subseteq C$ 可测（零测）但不是 Borel 集——否则 $W = H(H^{-1}(W))$ 将是 Borel 集。

\secn{2.12 习题选编}

\begin{itemize}

\item 证明：可数集的外测度为零；外测度为零的集合必可测。

\item 证明：$E$ 可测 $\iff$ 存在开集列从外部逼近（正则性演练）。

\item （Smith–Volterra–Cantor 集）每次删除区间长 $\frac{1}{4^n}$ 量级的中间开区间，得到正测度的无处稠密完全集——"胖 Cantor 集"。

\item 证明：$\mathbb{R}$ 上平移不变、规范化的 Borel 测度必为 Lebesgue 测度（常数倍）。

\item （Lebesgue–Stieltjes）给定单调递增右连续 $F$，写出 $\mu_F$ 在生成代数上的取值并验证预测度性质。

\item 常见错误辨析：$m^*$ 只在\textbf{可测集}上可数可加；次可加性对一般集合成立。

\end{itemize}

\begin{quote}
\textbf{本章小结}：外测度（覆盖定义、次可加、平移不变）$\to$ Carathéodory 可测性（$\sigma$-代数 + 完备测度 + 可数可加）$\to$ 结构定理（开集/F_$\sigma$ 逼近，Borel + 零测分解）$\to$ Vitali 不可测集 $\to$ 扩张定理（一般测度、L–S 测度）$\to$ 完备化与乘积测度预备。
\end{quote}

\chap{第三章　可测函数及其性质}

本章定义可测函数并研究其运算封闭性（四则、上下确界、极限）；建立三大逼近/收敛定理——简单函数逼近、Egorov 定理（几乎一致收敛）、Lusin 定理（近似连续），最后引入依测度收敛并证明 Lebesgue 定理与 Riesz 定理。

\secn{3.1 可测函数的定义}

设 $E \subseteq \mathbb{R}^n$ 可测，$f$ 是 $E$ 上的广义实值函数。

\begin{definition}[Lebesgue 可测函数]
若对任意 $t \in \mathbb{R}$，

\[
E\{f > t\} = \{x \in E : f(x) > t\} \in \mathcal{M},
\]

则称 $f$ 是 $E$ 上的（Lebesgue）可测函数。
\end{definition}

\begin{proposition}[等价刻画]
$f$ 可测 $\iff$ 以下之一成立：
\end{proposition}

（1）$\forall t$，$E\{f \ge t\} \in \mathcal{M}$；

（2）$\forall t$，$E\{f < t\} \in \mathcal{M}$；

（3）$\forall t$，$E\{f \le t\} \in \mathcal{M}$。

\begin{proof}
例如（2）$\Rightarrow$（3）：$E\{f \le t\} = E \setminus E\{f > t\}$ 由可测集运算封闭；（3）$\Rightarrow$（2）：$E\{f < t\} = \bigcup_k E\{f \le t - \tfrac{1}{k}\}$。
\end{proof}

\begin{remark}
连续函数可测（$\{f > t\}$ 是开集与 $E$ 的交）；示性函数 $\chi_A$ 可测 $\iff$ $A$ 可测。
\end{remark}

\secn{3.2 可测函数的运算}

\begin{proposition}[逐点运算的封闭性]
设 $f$、$g$ 在 $E$ 上可测：
\end{proposition}

（1）$f + c$ 可测（$E\{f + c > t\} = E\{f > t - c\}$）；

（2）$c f$ 可测（$c > 0$ 时 $E\{cf > t\} = E\{f > t/c\}$）；

（3）$\vert f \vert$ 可测（$E\{\vert f\vert > t\} = E\{f > t\} \cup E\{f < -t\}$）；

（4）$\vert f \vert^{\alpha}$ 可测（$\alpha > 0$；$E\{\vert f\vert^{\alpha} > t\} = E\{\vert f\vert > t^{1/\alpha}\}$，$t > 0$）；

（5）$\max(f, g)$、$\min(f, g)$ 可测（$E\{\max(f,g) > t\} = E\{f > t\} \cup E\{g > t\}$ 等）；

（6）可测函数的和、差、积、商（分母非零）可测（先证 $f^2$、$fg = \frac{1}{4}[(f+g)^2 - (f-g)^2]$ 型恒等式）。

\begin{proposition}[列的运算]
若 $\{f_n\}$ 可测，则

\[
E\left\{\sup_n f_n > t\right\} = \bigcup_n E\{f_n > t\}, \qquad E\left\{\inf_n f_n \ge t\right\} = \bigcap_n E\{f_n \ge t\},
\]

故 $\sup_n f_n$、$\inf_n f_n$、$\varlimsup_n f_n = \inf_n \sup_{k \ge n} f_k$、$\varliminf_n f_n = \sup_n \inf_{k \ge n} f_k$ 均可测。
\end{proposition}

\begin{example}[示性函数]
$A \subseteq E$：$\chi_A$ 可测 $\iff$ $A$ 可测。
\end{example}

\begin{example}[Dirichlet 函数的可测性]
$D(x) = \chi_{\mathbb{Q} \cap [0,1]}$ 可测（可数集是零测集），且

\[
D(x) = \lim_{m \to \infty} \lim_{n \to \infty} \left(\cos(m! \pi x)\right)^{2n},
\]

即 $D$ 是连续函数列的（重）极限——可测函数类对极限运算封闭的体现。
\end{example}

\secn{3.3 复合与 Borel 可测函数}

\begin{definition}[Borel 可测]
$f: E \to \mathbb{R}$ 称为 Borel 可测，若 $\forall t$，$\{f > t\}$ 是 Borel 集。
\end{definition}

\begin{proposition}[连续复合]
$f$ Lebesgue 可测、$g$ 连续 $\Rightarrow$ $g \circ f$ Lebesgue 可测。
\end{proposition}

\begin{proof}
要证 $\{x : g(f(x)) > t\}$ 可测。由 $g$ 连续，$\{y : g(y) > t\}$ 是开集，写成互不相交开区间之并 $\bigcup_j I_j$，于是

\[
\{x : g(f(x)) > t\} = \bigcup_j \{x : f(x) \in I_j\},
\]

每个 $\{f \in I_j\} = E\{f < b_j\} \cap E\{f > a_j\}$ 可测。
\end{proof}

\begin{remark}
$g$ 仅 Lebesgue 可测时结论可以失败（$f$ 连续、$g$ Lebesgue 可测，$g \circ f$ 可能不可测）；但 $f$ Borel 可测、$g$ Lebesgue 可测时 $g \circ f$ 可测。
\end{remark}

\begin{proposition}[Borel 逼近]
若 $f$ 是 Lebesgue 可测，则存在 Borel 可测函数 $g$ 使 $f = g$ a.e.。
\end{proposition}

\begin{proof}
Step 1：示性函数——$A$ 可测时存在 Borel 集 $B$ 使 $m(A \triangle B) = 0$（第二章结构定理），则 $\chi_A = \chi_B$ a.e.。Step 2：简单函数——$\varphi = \sum_j c_j \chi_{E_j}$，取 Borel 的 $B_j$ 替换 $E_j$ 得 $\psi$，$\varphi = \psi$ a.e.。Step 3：一般 $f$ 取简单函数逼近列 $\varphi_n \to f$，对应的 Borel 简单列 $\psi_n$，令 $g = \varlimsup_n \psi_n$ 即可。
\end{proof}

\secn{3.4 简单函数逼近}

\begin{definition}[简单函数]
$\varphi = \sum_{j=1}^{n} c_j \chi_{E_j}$，其中 $E_j$ 互不相交的可测集、$c_j$ 互不相同。
\end{definition}

\begin{theorem}[3.1（简单函数逼近）]
设 $f \ge 0$ 可测，则存在简单函数列 $\varphi_n \ge 0$ 使 $\varphi_n \uparrow f$（逐点单调递增收敛），且：
\end{theorem}

（i）若 $f$ 有界，则 $\varphi_n \rightrightarrows f$（一致收敛）；

（ii）若 $f$ 无界，则 $\varphi_n$ 仍处处收敛于 $f$。

\begin{proof}
二进制切割：

\[
E_{j,n} = \left\{ \frac{j}{2^n} \le f < \frac{j+1}{2^n} \right\},\quad j = 0, 1, \ldots, n2^n - 1; \qquad E_{n2^n, n} = \{f \ge n\},
\]

\[
\varphi_n = \sum_{j=0}^{n 2^n - 1} \frac{j}{2^n} \chi_{E_{j,n}} + n\,\chi_{\{f \ge n\}}.
\]

则在 $\{f < n\}$ 上 $\varphi_n \le f < \varphi_n + \tfrac{1}{2^n}$；$\varphi_n$ 简单、$\varphi_n \le f$、$\varphi_{n+1} \ge \varphi_n$，故 $\varphi_n \uparrow f$；$f$ 有界时一致。
\end{proof}

\begin{remark}
一般 $f$ 用 $f = f^+ - f^-$ 分解为两列递增简单函数之差。
\end{remark}

\secn{3.5 Egorov 定理}

\begin{theorem}[3.2（Egorov）]
设 $m(E) < \infty$，$\{f_n\}$ 是 $E$ 上 a.e. 有限的可测函数列，且 $f_n \to f$ a.e. 于 $E$（$f$ 也 a.e. 有限），则 $\forall \delta > 0$，存在可测子集 $E_\delta \subset E$，$m(E_\delta) < \delta$，使 $\{f_n\}$ 在 $E \setminus E_\delta$ 上一致收敛于 $f$（"几乎一致收敛"）。
\end{theorem}

\begin{proof}
由 a.e. 收敛，$m\{x \in E : f_n(x) \not\to f(x)\} = 0$，不妨设处处收敛。对每个 $k$，集合

\[
A_{N,k} = \bigcap_{n \ge N}\left\{x \in E : \vert f_n(x) - f(x)\vert < \tfrac{1}{k}\right\}
\]

随 $N$ 递增，且 $\bigcup_N A_{N,k} \supseteq E \setminus \{f_n \not\to f\}$；由测度下连续性（$m(E) < \infty$），$\lim_N m(E \setminus A_{N,k}) = 0$，故可取 $N_k$ 使 $m(E \setminus A_{N_k, k}) < \delta / 2^k$。

令 $E_\delta = \bigcup_{k=1}^{\infty} (E \setminus A_{N_k,k})$，则 $m(E_\delta) < \delta$；且 $\forall x \in E \setminus E_\delta$，$\forall k$，$\forall n \ge N_k$ 有 $\vert f_n(x) - f(x)\vert < \tfrac{1}{k}$，即一致收敛。
\end{proof}

\begin{remark}
条件 $m(E) < \infty$ 不可去掉。
\end{remark}

\begin{example}[$m(E) = \infty$ 时失效]
$f_n = \chi_{[n, n+1]}$ 于 $\mathbb{R}$：$f_n \to 0$ a.e.，但对任何 $\delta$，$\mathbb{R} \setminus E_\delta$ 含任意长的区间，$f_n$ 在其上不一致收敛于 0。
\end{example}

\begin{example}[滑脊函数]
$f_n$ 依次取高 $1$、底长趋于 $0$ 且位置向 $0$ 滑动的"脊"（$f_n = \chi_{(\frac{1}{n+1}, \frac{1}{n}]}$ 型）：$f_n \to 0$ a.e.，但不存在子列在任何包含 $0$ 的区间上一致收敛。
\end{example}

\secn{3.6 Lusin 定理}

\begin{theorem}[3.3（Lusin）]
设 $E$ 可测，$f$ 是 $E$ 上 a.e. 有限的可测函数，则 $\forall \delta > 0$，存在闭集 $F_\delta \subseteq E$，$m(E \setminus F_\delta) < \delta$，使 $f$ 在 $F_\delta$ 上连续。
\end{theorem}

\begin{proof}
Step 1（简单函数）：$f = \sum_{j=1}^n c_j \chi_{E_j}$（$E_j$ 互不相交可测）。由正则性取闭集 $F_j \subseteq E_j$ 使 $m(E_j \setminus F_j) < \delta/n$，令 $F = \bigcup_j F_j$。$F$ 是有限个互不相交闭集之并，$f$ 在每个 $F_j$ 上取常值，故在 $F$ 上连续（不相交闭集上可分别连续的函数在并上连续）。

Step 2（有界函数）：由逼近定理取有界一致收敛的简单列 $\varphi_n \rightrightarrows f$（截断后可设有界）。对每个 $\varphi_n$ 取 Step 1 的闭集 $F_n$（$m(E \setminus F_n) < \delta/2^n$），令 $F^* = \bigcap_n F_n$（闭集），则 $m(E \setminus F^*) < \delta$，且在 $F^*$ 上 $\varphi_n \rightrightarrows f$；连续函数（在闭集上）的一致极限连续。

Step 3（一般情形）：截断 $f^{(N)} = \max(-N, \min(f, N))$ 有界可测，$\{f \ne f^{(N)}\} \downarrow$ 零测集；取 $N$ 足够大后对 $f^{(N)}$ 用 Step 2，再并入一个小测度集。
\end{proof}

\begin{quote}
\textbf{Lusin 定理的含义}：可测函数"差不多"是连续函数——差别集中在任意小测度的集合上。它与 Egorov 定理一起构成"可测 ⟹ 近似好"的图景。
\end{quote}

\secn{3.7 依测度收敛}

\begin{definition}[依测度收敛]
设 $f$、$f_n$ 是 $E$ 上 a.e. 有限的可测函数。若 $\forall \sigma > 0$，

\[
m\{x \in E : \vert f_n(x) - f(x)\vert \ge \sigma\} \to 0 \quad (n \to \infty),
\]

则称 $f_n$ \textbf{依测度收敛}（converge in measure / in probability）于 $f$，记 $f_n \Rightarrow f$。
\end{definition}

\begin{theorem}[3.4（Lebesgue）]
设 $m(E) < \infty$，$f_n \to f$ a.e. 于 $E$，则 $f_n \Rightarrow f$。
\end{theorem}

\begin{proof}
由 Egorov 定理，$\forall \delta > 0$ 取 $E_\delta$ 使 $m(E_\delta) < \delta$ 且在 $E \setminus E_\delta$ 上 $f_n \rightrightarrows f$。当 $n$ 充分大时

\[
\{\vert f_n - f\vert \ge \sigma\} \subseteq E_\delta,
\]

故 $m\{\vert f_n - f\vert \ge \sigma\} < \delta$；由 $\delta$ 任意得证。
\end{proof}

\begin{remark}
逆命题不成立（滑脊例：$f_n \Rightarrow 0$ 但处处不收敛）；$m(E) = \infty$ 时 a.e. 收敛也不蕴含依测度收敛（$\chi_{[n,n+1]}$ 例）。
\end{remark}

\secn{3.8 Riesz 定理}

\begin{theorem}[3.5（Riesz）]
若 $f_n \Rightarrow f$（依测度收敛），则存在子列 $\{f_{n_k}\}$ 使 $f_{n_k} \to f$ a.e.
\end{theorem}

\begin{proof}
对每个 $k$，由依测度收敛取 $n_k$（$n_1 < n_2 < \cdots$）使

\[
m E_k := m\{\vert f_{n_k} - f\vert \ge 1/k\} < 2^{-k}.
\]

则 $\sum_k m E_k < \infty$，对上极限集用测度的次可加性与连续性：

\[
m\left(\varlimsup_{k \to \infty} E_k\right) = \lim_{N \to \infty} m\left(\bigcup_{k \ge N} E_k\right) \le \lim_{N \to \infty} \sum_{k \ge N} 2^{-k} = 0.
\]

对 $x \notin \varlimsup_k E_k$：存在 $K$，$k \ge K$ 时 $x \notin E_k$，即 $\vert f_{n_k}(x) - f(x)\vert < 1/k \to 0$，故 $f_{n_k}(x) \to f(x)$。
\end{proof}

\begin{quote}
\textbf{本章小结}：可测函数 = 用 $\{f > t\}$ 可测刻画的、对一切"好运算"（代数运算、上下确界、极限）封闭的函数类；简单函数逼近（构造积分的基础）；Egorov（a.e. 收敛 $\Rightarrow$ 几乎一致收敛，有限测度）；Lusin（可测 $\Rightarrow$ 近似连续）；依测度收敛与 a.e. 收敛经 Riesz 定理相互沟通——为第四章 Lebesgue 积分做好准备。
\end{quote}

\chap{第四章　Lebesgue 积分}

本章按"简单函数 $\to$ 非负可测函数 $\to$ 一般可积函数"三步建立 Lebesgue 积分，证明三大收敛定理（Levi 单调收敛、Fatou 引理、控制收敛定理），讨论积分与极限交换的大量典型例子、Riemann 可积性与 Lebesgue 可积性的关系，最后建立 Tonelli 与 Fubini 定理。

\secn{4.1 非负可测函数的积分}

\begin{definition}[非负可测函数空间]
$L^+ = \{f : f \text{ 在 } E \text{ 上非负可测}\}$（取值于 $[0, +\infty]$）。
\end{definition}

\begin{definition}[简单函数的积分]
对简单函数 $\varphi = \sum_{j=1}^{n} a_j \chi_{A_j}$（$A_j$ 两两不交可测），定义

\[
\int_E \varphi \, dm := \sum_{j=1}^n a_j \, m(A_j).
\]
\end{definition}

\begin{definition}[非负可测函数的积分]
对 $f \in L^+$，定义

\[
\int_E f \, dm := \sup\left\{ \int_E \varphi \, dm : 0 \le \varphi \le f,\ \varphi \text{ 简单可测} \right\}.
\]
\end{definition}

\begin{proposition}[基本性质]

\end{proposition}

（1）设 $0 \le \varphi \le f$，$\varphi$ 简单可测，则 $\int_E \varphi \, dm \le \int_E f \, dm$；

（2）（单调性）$0 \le f \le g$ $\Rightarrow$ $\int_E f \, dm \le \int_E g \, dm$；

（3）$m(E) = 0$ $\Rightarrow$ $\int_E f \, dm = 0$。

\begin{proposition}
设 $m(E) < \infty$，则 $\int_E f \, dm = 0 \iff f = 0$ a.e. 于 $E$。
\end{proposition}

\begin{proof}
"$\Leftarrow$"：设 $f = 0$ a.e.，对 $0 \le \varphi = \sum_j a_j \chi_{A_j} \le f$ 有 $m(A_j) = 0$，故 $\int_E \varphi = 0$，取上确界得 $\int_E f\,dm = 0$。

"$\Rightarrow$"：设 $m\{x \in E : f(x) > 0\} > 0$。由 $\{f > 0\} = \bigcup_k \{f > \tfrac{1}{k}\}$ 与测度下连续性，存在 $k$ 使 $m\{f > \tfrac{1}{k}\} > 0$。令 $\varphi = \tfrac{1}{k} \chi_{\{f > 1/k\}} \le f$，则

\[
\int_E f \, dm \ge \int_E \varphi \, dm = \tfrac{1}{k}\, m\{f > \tfrac{1}{k}\} > 0,
\]

矛盾。
\end{proof}

\begin{remark}
（1）（对集合的可加性）$E_1 \cap E_2 = \emptyset$ 时 $\int_{E_1 \cup E_2} f \, dm = \int_{E_1} f \, dm + \int_{E_2} f \, dm$（先对简单函数验证，再取 sup）。
\end{remark}

（2）（限制与示性函数）$A \subseteq E$ 可测时 $\int_A f \, dm = \int_E f \cdot \chi_A \, dm$：先对简单函数 $\varphi$ 验证 $\int_E \varphi \chi_A \, dm = \sum_j a_j\, m(A_j \cap A) = \int_A \varphi \, dm$，一般情形用递增简单列 $\varphi_j \nearrow f$（第三章逼近定理）+ 下节 MCT 取极限。

\begin{theorem}[线性性]
设 $f, g \in L^+$，$\alpha, \beta \ge 0$，则

\[
\int_E (\alpha f + \beta g)\, dm = \alpha \int_E f\, dm + \beta \int_E g\, dm.
\]
\end{theorem}

\begin{proof}
取简单可测 $\varphi_j \nearrow f$、$\psi_j \nearrow g$，则 $0 \le \alpha\varphi_j + \beta\psi_j \nearrow \alpha f + \beta g$；对简单函数线性性已证（直接按定义展开），再用下面的 MCT 取极限。
\end{proof}

\begin{theorem}[4.1（Levi 单调收敛定理，MCT）]
设 $\{f_j\} \subset L^+$，$f_j \le f_{j+1}$，$f_j \to f$ 逐点（从而 $f \in L^+$），则

\[
\int_E f \, dm = \lim_{j \to \infty} \int_E f_j \, dm.
\]
\end{theorem}

\begin{proof}
由单调性 $\int f_j$ 递增且 $\le \int f$，故 $\lim_j \int f_j \le \int f$。反之，任取简单 $0 \le \varphi \le f$ 与 $0 < t < 1$，令 $E_j = \{x : f_j(x) \ge t\,\varphi(x)\}$，则 $E_j$ 可测、$E_j \nearrow E$，且 $f_j \ge t\,\varphi\,\chi_{E_j}$。由 $\{f > t\varphi\}$ 的逼近与测度下连续性（$m(E)$ 可以取 $+\infty$，下连续性不需要有限）

\[
\lim_j \int_E f_j \, dm \ge \lim_j \int_E t\,\varphi\,\chi_{E_j}\,dm = t \int_E \varphi \, dm \xrightarrow{\,t \to 1\,} \int_E \varphi \, dm,
\]

对 $\varphi$ 取上确界得 $\lim_j \int f_j \ge \int f$。
\end{proof}

\begin{remark}[1]
单调性条件不可去掉：$f_j = j\,\chi_{(0, \frac{1}{j})} \to 0$ a.e.，但 $\int f_j = 1 \not\to 0$。
\end{remark}

\begin{remark}[2]
若把条件减弱为"$f_j \nearrow f$ a.e. 且某个 $\int_E f_{j_0} \, dm < +\infty$"，结论仍真。
\end{remark}

\begin{theorem}[4.2（逐项积分）]
设 $f_j \in L^+$，则

\[
\int_E \sum_{j=1}^{\infty} f_j \, dm = \sum_{j=1}^{\infty} \int_E f_j \, dm.
\]
\end{theorem}

\begin{proof}
令 $g_n = \sum_{j=1}^n f_j \nearrow \sum_{j=1}^\infty f_j$，由 MCT 与有限可加性

\[
\int_E \sum_{j=1}^{\infty} f_j \, dm = \lim_n \int_E g_n \, dm = \lim_n \sum_{j=1}^{n} \int_E f_j \, dm = \sum_{j=1}^{\infty} \int_E f_j \, dm. \; \]

（另一方向的说明：若 $\sum_j f_j$ 处处发散于 $+\infty$ 而 $m(E) > 0$，则积分也为 $+\infty$——取 $E_{j,k} = \{f_j \ge \tfrac1k\}$ 型集合控制即可。）
\end{proof}

\secn{4.2 积分与极限的交换}

先看一个反例（"sliding hump"，滑动的驼峰）：$f_j = \chi_{(j,\, j+1)}$，则 $f_j \to 0$ 逐点，$\int f_j = 1$，但 $\int \lim f_j = 0 < 1 = \lim \int f_j$——极限与积分交换需要条件。

\begin{theorem}[4.3（Fatou 引理）]
设 $\{f_j\} \subset L^+$，则

\[
\int_E \varliminf_{j \to \infty} f_j \, dm \le \varliminf_{j \to \infty} \int_E f_j \, dm.
\]
\end{theorem}

\begin{proof}
令 $g_k = \inf_{j \ge k} f_j \in L^+$，则 $g_k \nearrow \varliminf f_j$；对 $g_k$ 用 MCT：

\[
\int_E \varliminf_j f_j \, dm = \lim_k \int_E g_k \, dm = \sup_k \int_E g_k \, dm \le \varliminf_j \int_E f_j \, dm. \; \]
\end{proof}

\begin{remark}[1]
不等式可以严格（sliding hump 例）。
\end{remark}

\begin{remark}[2]
若 $f_j \to f$ a.e. 且 $f_j \le g$（$g$ 可积），对 $g - f_j \ge 0$ 用 Fatou：

\[
\int g - f \le \liminf_j \int (g - f_j) = \int g - \limsup_j \int f_j.
\]
\end{remark}

\begin{corollary}[4.4（控制收敛定理，DCT）]
设 $f_j \to f$ a.e. 于 $E$，$\vert f_j \vert \le g$（$g$ 可积），则 $f_j$、$f$ 均可积，且

\[
\lim_{j \to \infty} \int_E f_j \, dm = \int_E f \, dm.
\]
\end{corollary}

\begin{proof}
由 $\vert f_j\vert \le g$ 知 $f_j, f$ 可积。对 $g - f_j$ 与 $g + f_j$ 分别用注 2 的 Fatou 论证，得

\[
\int f \le \liminf \int f_j \le \limsup \int f_j \le \int f,
\]

故极限存在且相等。
\end{proof}

\begin{theorem}[4.5（一般的控制收敛定理）]
设 $f_j \to f$ a.e.，$\vert f_j\vert \le g_j$，$g_j \to g$ a.e.，$g$ 可积，且 $\lim_j \int_E g_j \, dm = \int_E g \, dm$，则 $f$ 可积且 $\lim_j \int f_j = \int f$。
\end{theorem}

\begin{proof}
$\vert f_j - f\vert \le g_j + g$，对非负函数 $g_j + g - \vert f_j - f\vert$ 用 Fatou：

\[
\int 2g - \limsup_j \int \vert f_j - f\vert \le \liminf_j \int (g_j + g - \vert f_j - f\vert) = \int 2g,
\]

即 $\limsup_j \int \vert f_j - f\vert \le 0$，故 $\int \vert f_j - f\vert \to 0$，而 $\left\vert \int f_j - \int f \right\vert \le \int \vert f_j - f\vert \to 0$。
\end{proof}

\begin{theorem}[积分的绝对连续性]
设 $f$ 在 $E$ 上可积，则 $\forall \varepsilon > 0$，$\exists \delta > 0$：只要 $A \subseteq E$ 可测且 $m(A) < \delta$，就有

\[
\int_A \vert f \vert \, dm < \varepsilon.
\]
\end{theorem}

\begin{proof}
先对 $f = \chi_E$（取 $\delta = \varepsilon$）与简单函数验证，一般情形用简单函数逼近 + MCT 控制 $\int \vert f - \varphi\vert < \varepsilon/2$ 后取 $\delta = \varepsilon/2$。
\end{proof}

\begin{example}[1（Beta 函数）]
$\displaystyle\int_0^1 x^{a-1}(1-x)^{b-1}\,dx = B(a,b)$（$a, b > 0$；非负可测，允许 $+\infty$ 值的逐项讨论）。
\end{example}

\begin{example}[2]
求 $\lim_n \int_0^n \left(1 - \frac{x}{n}\right)^n e^{x/2}\,dx$。
\end{example}

\begin{proof}[解]
令 $f_n = \left(1 - \frac{x}{n}\right)^n e^{x/2} \chi_{[0,n]}$。当 $x \ge 0$ 固定、$n > x$ 时，$\left(1 - \frac{x}{n}\right)^n \to e^{-x}$（且关于 $n$ 最终递增），故 $f_n \to e^{-x/2}$；又 $\left(1 - \frac{t}{n}\right)^n \le e^{-t}$（$0 \le t \le n$）给出控制 $\vert f_n\vert \le e^{-x/2}$ 可积。由 DCT：

\[
\lim_n \int_0^n \left(1-\frac{x}{n}\right)^n e^{x/2}\,dx = \int_0^\infty e^{-x/2}\,dx = 2.
\]
\end{proof}

\begin{example}[3]
$\displaystyle\int_0^1 \frac{\ln x}{1 - x}\,dx = -\frac{\pi^2}{6}$。
\end{example}

\begin{proof}[解]
$\frac{1}{1-x} = \sum_{n=0}^{\infty} x^n$（$0 \le x < 1$），$\ln x < 0$，由逐项积分定理（非负情形对 $-\frac{\ln x}{1-x} = \sum_n (-\ln x)\, x^n$ 应用）

\[
-\int_0^1 \frac{\ln x}{1 - x}\,dx = \sum_{n=0}^{\infty} \int_0^1 (-\ln x)\, x^n \,dx = \sum_{n=0}^{\infty} \frac{1}{(n+1)^2} = \frac{\pi^2}{6}. \; \]
\end{proof}

\begin{example}[4（Riemann 可积 ⟹ Lebesgue 可积）]
设 $f$ 在 $[a,b]$ 上 Riemann 可积，则 $f$ Lebesgue 可积，且两个积分值相等。
\end{example}

\begin{proof}
取分划列 $P_n: a = x_0^n < \cdots < x_{N_n}^n = b$，$\Vert P_n\Vert \to 0$（例如等分）。令

\[
u_n = \sum_k m_k^n \,\chi_{(x_{k-1}^n,\, x_k^n]} \le f \le v_n = \sum_k M_k^n \,\chi_{(x_{k-1}^n,\, x_k^n]},
\]

其中 $m_k^n, M_k^n$ 为 $f$ 在子区间上的下、上确界（Riemann 可积函数有界、a.e. 连续，故可测）。$u_n$ 的积分是 Darboux 下和 $\underline{S}_n$、$v_n$ 是上和 $\overline{S}_n$，由 Riemann 可积性 $\underline{S}_n, \overline{S}_n \to \int_a^b f\,dx$。取一列加细分划使 $\underline{S}_n \nearrow$，由 MCT 与夹逼：

\[
(L)\int_a^b f \, dm = \lim_n \int u_n = \lim_n \underline{S}_n = \int_a^b f \, dx. \; \]
\end{proof}

\begin{example}[5]
$\lim_n \int_0^1 \frac{n^3 x^2\, e^{-n^2 x}}{1 + x}\,dx = 0$。
\end{example}

\begin{proof}[解]
$f_n(x) = \frac{n^3 x^2}{1+x} e^{-n^2 x} \to 0$（$x > 0$；$f_n(0)=0$）。求导分析最大值可得 $\vert f_n(x)\vert \le \frac{C}{\sqrt{x}}$（$C$ 与 $n$ 无关），而 $x^{-1/2}$ 在 $(0,1)$ 可积，由 DCT 极限为 $0$。
\end{proof}

\begin{example}[6]
（a）$\lim_n \int_0^1 \frac{n\sqrt{x}}{1 + n^2 x^2}\,dx = 0$；（b）$\displaystyle\int_0^1 \frac{x^{\alpha} - 1}{\ln x}\,dx = \ln(\alpha + 1)$（$\alpha > -1$）。
\end{example}

\begin{proof}[解]
（a）$f_n \to 0$（$x > 0$）。用 $n x \le \frac{1}{2}(n^2 x^2 + 1)$（即 $2nx \le n^2x^2 + 1$）：

\[
\frac{n\sqrt{x}}{1 + n^2 x^2} \le \frac{n\sqrt{x}}{2 n x} = \frac{1}{2}\,x^{-1/2} \in L^1(0,1),
\]

由 DCT 极限为 $0$。

（b）$\frac{x^{\alpha} - 1}{\ln x} = \int_0^{\alpha} x^t \,dt$，由 Fubini/Tonelli 交换积分次序：

\[
\int_0^1 \frac{x^{\alpha} - 1}{\ln x}\,dx = \int_0^\alpha \left(\int_0^1 x^t\,dx\right) dt = \int_0^\alpha \frac{dt}{t+1} = \ln(\alpha + 1). \; \]
\end{proof}

\begin{example}[7]
$\displaystyle\int_0^1 \!\!\int_0^1 \frac{dx\,dy}{1 - xy} = \frac{\pi^2}{6}$。
\end{example}

\begin{proof}[解]
$f(x,y) = \frac{1}{1-xy} \ge 0$ 可测，由 $\frac{1}{1-t} = \sum_n t^n$（$0 \le t < 1$）与 Tonelli 逐项积分：

\[
\int_0^1\!\!\int_0^1 \frac{dx\,dy}{1-xy} = \sum_{n=0}^\infty \int_0^1\!\!\int_0^1 x^n y^n \,dx\,dy = \sum_{n=0}^\infty \frac{1}{(n+1)^2} = \frac{\pi^2}{6}. \; \]
\end{proof}

\begin{example}[交换积分次序不成立的例子]
令

\[
f(x, y) = \frac{x^2 - y^2}{(x^2 + y^2)^2}, \qquad (x, y) \in (0,1)\times(0,1).
\]

则 $\int_0^1 f\,dy = \left[\frac{y}{x^2 + y^2}\right]_{y=0}^{y=1} = \frac{1}{1 + x^2}$，故一个次序的累次积分为 $\int_0^1 \frac{dx}{1+x^2} = \frac{\pi}{4}$；而另一次序（对 $x$ 先积分）得 $-\frac{\pi}{4}$。两个累次积分不相等——$f \notin L^1((0,1)^2)$，Fubini 定理的条件不可省略。
\end{example}

\secn{4.3 Fubini 定理}

\begin{theorem}[4.6（Tonelli）]
设 $(X, \mathcal{X}, \mu)$、$(Y, \mathcal{Y}, \nu)$ 为 $\sigma$-有限测度空间，$f$ 为 $X \times Y$ 上的非负可测函数，则

\[
\int_{X \times Y} f \,d(\mu \times \nu) = \int_X \left[ \int_Y f \,d\nu \right] d\mu = \int_Y \left[ \int_X f \,d\mu \right] d\nu.
\]
\end{theorem}

（三个量都允许为 $+\infty$。）

\begin{proof}
分三步：

（1）示性函数 $f = \chi_{A \times B}$（$A \in \mathcal{X}$，$B \in \mathcal{Y}$）：三边都等于 $\mu(A)\,\nu(B)$；由线性性推广到简单函数。

（2）单调类方法：令 $\mathcal{A} = \{E \in \mathcal{X} \otimes \mathcal{Y} : \text{结论对 } \chi_E \text{ 成立}\}$，由（1）$\mathcal{A}$ 包含矩形全体生成的代数，且对递增极限封闭，由单调类定理 $\mathcal{A} = \mathcal{X} \otimes \mathcal{Y}$。

（3）一般非负可测 $f$：取简单函数 $f_j \nearrow f$，对 $f_j$ 结论成立，由 MCT 两次取极限传给 $f$。
\end{proof}

\begin{theorem}[4.7（Fubini）]
设 $(X, \mathcal{X}, \mu)$、$(Y, \mathcal{Y}, \nu)$ $\sigma$-有限，$f \in L^1(\mu \times \nu)$，则

\[
\int_{X\times Y} f \,d(\mu \times \nu) = \int_X\left[\int_Y f\,d\nu\right]d\mu = \int_Y\left[\int_X f\,d\mu\right]d\nu,
\]
\end{theorem}

且对 a.e. 的 $x$，内层积分 $\int_Y f(x, \cdot)\,d\nu$ 存在（有限）。

\begin{proof}
分解 $f = f^+ - f^-$，对 $f^\pm$（非负可测）用 Tonelli：

\[
\int_X \left[\int_Y f^+ \,d\nu\right] d\mu = \int_{X\times Y} f^+ \,d(\mu\times\nu) < \infty,
\]

故 $\int_Y f^+\,d\nu < \infty$ 对 a.e. $x$ 成立，同理 $f^-$；对这些 $x$ 可逐项相减：

\[
\int_X \left[\int_Y f \,d\nu\right] d\mu = \int_X \left[\int_Y f^+ \,d\nu\right]d\mu - \int_X\left[\int_Y f^- \,d\nu\right]d\mu = \int f^+ - \int f^- = \int f. \; \]
\end{proof}

\begin{remark}[条件不可削弱]
若仅有两个累次积分存在且相等，未必有 $f \in L^1$、重积分存在（见 4.2 末的反例）；正确的使用方式是：\textbf{非负用 Tonelli，可积性未知时先用 Tonelli 验证 $\int\int \vert f\vert < \infty$，再用 Fubini}。
\end{remark}

\begin{quote}
\textbf{本章小结}：三步构造（简单 $\to$ 非负 $\to$ 一般）$\to$ 三大收敛定理（MCT、Fatou、DCT 及一般形式）$\to$ 积分绝对连续性 $\to$ 典型极限积分问题（$(1-x/n)^n$ 型、参数积分、$\pi^2/6$ 型级数）$\to$ Riemann 与 Lebesgue 可积性的关系 $\to$ Tonelli/Fubini（非负随便换序、可积才能换序）。
\end{quote}

\chap{第五章　微分与不定积分}

本章回答第一章提出的问题：$F(x) = \int_a^x f\,dm$ 何时有 $F' = f$、Newton–Leibniz 公式成立的最小条件是什么。主线：单调函数几乎处处可微（Lebesgue）$\to$ 有界变差函数（Jordan 分解）$\to$ 绝对连续函数（N–L 公式成立的充要条件）$\to$ Lebesgue 微分定理（Hardy–Littlewood 极大函数与 Vitali 覆盖引理）$\to$ 换元、分部积分与 Stieltjes 积分。

\secn{5.1 单调函数的可微性}

\begin{definition}[Dini 导数]
\[
D^+ F(x) = \varlimsup_{h \to 0^+} \frac{F(x+h)-F(x)}{h}, \qquad D_+ F(x) = \varliminf_{h \to 0^+} \frac{F(x+h)-F(x)}{h},
\]
\end{definition}

同理用左极限定义 $D^- F$、$D_- F$（共四个 Dini 导数；$F$ 可微 $\iff$ 四者相等且有限）。

\begin{lemma}[Rising Sun 引理 / Riesz 引理]
设 $F \in C[a, b]$，则点集 $\{x \in (a,b) : \exists\, y > x,\ F(y) > F(x)\}$ 是开集；若记其构成区间为 $(a_k, b_k)$，则 $F(b_k) \ge F(a_k)$。
\end{lemma}

\begin{theorem}[5.1（Lebesgue）]
单调函数几乎处处可微；导数 $F'$ 可积，且

\[
\int_a^b F'\,dm \le F(b) - F(a).
\]
\end{theorem}

\begin{proof}[证明思路]
Step 1：Rising Sun 引理（反证 + 最值点处的导数符号矛盾）。Step 2：对 $F(x) + \varepsilon x$ 与 $-F$ 应用，得 $D^+ F \ge 0$ a.e. 与对称不等式。Step 3：对例外集 $\{D^+F > 0\} \cap \{D_-F < 0\}$ 型集合用 Vitali 覆盖引理（5.8）控制测度为零。Step 4：$\int F' \le F(b) - F(a)$ 由 Fatou 引理对差商取极限得到。
\end{proof}

\begin{example}[Cantor 函数]
Cantor 函数 $G$：$G' = 0$ a.e. 但 $G(1) - G(0) = 1 \ne 0 = \int_0^1 G'\,dm$——不等号可以严格。
\end{example}

\secn{5.2 有界变差函数}

\begin{definition}[全变差]
\[
V_a^b(F) = \sup_{P} \sum_{j=1}^{n} \vert F(x_j) - F(x_{j-1}) \vert,
\]
\end{definition}

其中 $P: a = x_0 < x_1 < \cdots < x_n = b$ 取遍 $[a,b]$ 的分割。$V_a^b(F) < \infty$ 的函数称为\textbf{有界变差函数}，记 $BV[a,b]$。

\begin{example}
单调函数 $\in BV$，且 $V_a^b = \vert F(b) - F(a)\vert$；Lipschitz 函数 $\in BV$。
\end{example}

\begin{example}[连续但非 BV]
$F(x) = x \sin(1/x)$ 于 $[0,1]$：在极值点处的振幅之和发散，$V_0^1(F) = +\infty$。
\end{example}

\begin{theorem}[5.2（Jordan 分解）]
$F \in BV[a,b] \iff$ 存在递增函数 $g, h$ 使 $F = g - h$。
\end{theorem}

\begin{proof}
取变差函数 $g(x) = V_a^x(F)$（递增），令 $h = g - F$；对 $x < y$，

\[
h(y) - h(x) = V_x^y(F) - (F(y) - F(x)) \ge 0,
\]

故 $h$ 递增。
\end{proof}

\begin{corollary}
BV 函数几乎处处可微（由 Lebesgue 定理）；BV 函数的不连续点至多可数且均为第一类不连续点。
\end{corollary}

\begin{proposition}
$\int_a^b \vert F'\vert \,dm \le V_a^b(F)$（等号未必成立：Cantor 函数 $V = 1$ 但 $\int \vert G'\vert = 0$）。
\end{proposition}

\secn{5.3 绝对连续函数}

\begin{definition}[绝对连续，AC]
$\forall \varepsilon > 0$，$\exists \delta > 0$：对任意有限个互不相交的区间 $(a_k, b_k)$，

\[
\sum_k (b_k - a_k) < \delta \implies \sum_k \vert F(b_k) - F(a_k) \vert < \varepsilon.
\]
\end{definition}

\begin{example}
Lipschitz $\Rightarrow$ AC；不定积分 $F(x) = \int_a^x f\,dm$（$f$ 可积）是 AC 的——这正是积分的绝对连续性（第四章）。
\end{example}

\begin{proposition}
AC $\Rightarrow$ 一致连续、连续、BV（AC 函数是连续 BV 函数）。
\end{proposition}

\textbf{反例（Cantor 函数）.} $G$ 连续、递增（从而 BV），但非 AC：它把零测的 Cantor 集映为正测度集 $[0,1]$（AC 函数把零测集映为零测集）。

\begin{remark}[真包含链]
\[
\{\text{Lipschitz}\} \subsetneq AC \subsetneq BV \cap C \subsetneq BV \subsetneq \{\text{递增差}\}.
\]
\end{remark}

\secn{5.4 Newton–Leibniz 公式}

\begin{theorem}[5.3（N–L 公式）]
$F$ 在 $[a,b]$ 上绝对连续 $\iff$ $F'$ 可积且

\[
F(x) - F(a) = \int_a^x F'(t)\,dt, \qquad \forall\, x \in [a,b].
\]
\end{theorem}

\begin{proof}
（$\Leftarrow$）不定积分自动 AC（积分绝对连续性）。

（$\Rightarrow$）分两步。

（1）\textbf{引理}：AC 且 $F' = 0$ a.e. $\Rightarrow$ $F$ 为常数。反证：若 $F(x_2) > F(x_1)$（$x_1 < x_2$），取例外零测集的小开覆盖（总长 $< \delta$），在覆盖之外的每个点 $F' = 0$，可用 Rising Sun/覆盖论证构造有限个不相交区间使 $\sum \vert F(b_k) - F(a_k)\vert < \varepsilon$ 但总增量 $\ge F(x_2) - F(x_1)$，与 AC 矛盾。

（2）一般情形：令 $G(x) = \int_a^x F'\,dm$，则 $G$ AC 且 $G' = F'$ a.e.（由下一节的微分定理，或对 AC 函数直接验证），故 $F - G$ AC 且 $(F-G)' = 0$ a.e.，由（1）$F - G$ 为常数，取 $x = a$ 定出常数。
\end{proof}

\begin{quote}
Newton–Leibniz 公式成立的函数类\textbf{恰好}是绝对连续函数——第一章问题（2）至此完全解决；Cantor 函数说明 BV/连续都不够。
\end{quote}

\secn{5.5 Lebesgue 微分定理}

\begin{theorem}[5.4（Lebesgue 微分定理）]
设 $f \in L^1(\mathbb{R}^n)$，则

\[
\lim_{r \to 0} \frac{1}{m(B(x,r))} \int_{B(x,r)} f \, dm = f(x) \qquad \text{a.e. } x,
\]

即使得等式成立的 $x$ 全体（Lebesgue 点）的余集是零测集。
\end{theorem}

\begin{proof}[证明思路]
定义中心极大函数

\[
Mf(x) = \sup_{r > 0} \frac{1}{m(B(x,r))} \int_{B(x,r)} \vert f \vert \, dm .
\]

（1）弱 (1,1) 不等式：$m\{Mf > \lambda\} \le \frac{C}{\lambda} \Vert f\Vert_1$——用 Vitali 覆盖引理（5.8）的 5r-型覆盖论证；

（2）$M$ 显然强 $(\infty, \infty)$；插值得弱/强 $(p, p)$（$p > 1$）；

（3）对稠密集（$C_c$ 连续函数，结论逐点显然）逼近 + 极大函数控制余集：

\[
\Omega = \left\{ \varlimsup_{r \to 0} \left\vert \frac{1}{m(B_r)}\int_{B(x,r)} f - f(x) \right\vert > 0 \right\} \subseteq \{M(f - \varphi) > \tfrac{\varepsilon}{3}\} \cup \{\vert f - \varphi\vert > \tfrac{\varepsilon}{3}\},
\]

令 $\varepsilon \downarrow 0$ 得 $m(\Omega) = 0$。
\end{proof}

\begin{corollary}[积分的微分定理]
$F(x) = \int_a^x f\,dm$（$f \in L^1$）在 Lebesgue 点处可微且 $F' = f$；一般测度 $\mu \ll m$ 时 $\frac{\mu(B(x,r))}{m(B(x,r))} \to \frac{d\mu}{dm}$ a.e.（Radon–Nikodym 导数的实现）。
\end{corollary}

\begin{remark}
有界变差函数的导数可积，但 $\int_a^b F' \le F(b) - F(a)$ 可严格（Cantor 函数）；等号成立 $\iff$ $F$ AC。
\end{remark}

\secn{5.6 换元积分法与分部积分}

\begin{theorem}[5.5（换元法）]
设 $\varphi$ 在 $[a,b]$ 上绝对连续、严格递增，$f$ 可积，则

\[
\int_{\varphi(a)}^{\varphi(b)} f(y)\,dy = \int_a^b f(\varphi(x))\,\varphi'(x)\,dx .
\]
\end{theorem}

\begin{theorem}[5.6（分部积分）]
设 $F, G \in AC[a,b]$，则

\[
\int_a^b F\,G'\,dm = F(b)G(b) - F(a)G(a) - \int_a^b F'\,G\,dm .
\]
\end{theorem}

\begin{proof}
两式都由 N–L 公式直接验证：换元法对 $f = \chi_{(c,d]}$ 即 5.3 中"AC 函数把零测集映为零测集"的测度等式，再用典型方法逼近；分部积分对 $FG$ 用乘积法则（$FG \in AC$，$(FG)' = F'G + FG'$ a.e.）。
\end{proof}

\secn{5.7 Stieltjes 积分}

\begin{definition}[Riemann–Stieltjes 积分]
对 $[a,b]$ 上的 $f$ 与递增 $\alpha$，定义

\[
\int_a^b f \,d\alpha = \lim_{\Vert P\Vert \to 0} \sum_j f(\xi_j)\left[\alpha(x_j) - \alpha(x_{j-1})\right]
\]
\end{definition}

（当极限存在时）。

\begin{proposition}
$\alpha$ 递增右连续时，R–S 积分与 Lebesgue–Stieltjes 积分一致：

\[
\int_a^b f \,d\alpha = \int_{(a,b]} f \,d\mu_\alpha .
\]
\end{proposition}

\begin{example}[跳跃函数]
$\alpha$ 为阶梯函数（跳跃 $\Delta\alpha_j$ 在 $x_j$ 处）时，积分化为级数 $\sum_j f(x_j)\,\Delta\alpha_j$。
\end{example}

\begin{remark}[概率论]
$EX = \int x \,dF(x)$（分布函数的 R–S 积分）；一般期望理论就是 L–S 积分。
\end{remark}

\secn{5.8 Vitali 覆盖引理}

\begin{lemma}[Vitali 覆盖引理]
设 $\mathcal{V}$ 是 $E \subseteq \mathbb{R}^n$ 的 Vitali 覆盖：$\forall x \in E$、$\forall \varepsilon > 0$，存在 $I \in \mathcal{V}$ 使 $x \in I$ 且 $\mathrm{diam}(I) < \varepsilon$。则可选出可数个两两不交的 $\{I_k\} \subset \mathcal{V}$，使

\[
m\left(E \setminus \bigcup_k I_k\right) = 0.
\]
\end{lemma}

\begin{proof}[证明思路]
贪心选取：取有界区域覆盖，第 $k$ 步在未被覆盖部分选取直径尽可能大的区间；对余下部分用 5r 估计（每个未选中的区间被某个选中的、至多 5 倍直径的区间覆盖），余集测度趋于零。
\end{proof}

\begin{remark}
Vitali 覆盖引理是 H–L 极大函数弱 (1,1) 与微分定理的标准工具；与 2.3 的 Vitali 不可测集同源。
\end{remark}

\secn{5.9 例题与反例汇总}

\begin{example}[1（Cantor 函数）]
连续、递增、$G' = 0$ a.e.、$G(1) > G(0)$：N–L 失败（非 AC）；同时说明"$F' = 0$ a.e. $\Rightarrow$ 常数"需要 AC。
\end{example}

\begin{example}[2（$F'$ 存在但不可积）]
$F(x) = x^2 \sin(x^{-2})$ 于 $[0,1]$（$F(0)=0$）：处处可微，但 $F'(x) \sim 2\cos(x^{-2})/x$ 在 $0$ 附近不可积，$F' \notin L^1$。
\end{example}

\begin{example}[3（跳跃函数）]
阶梯函数 BV、不 AC；一般 BV 函数有唯一分解

\[
F = F_{AC} + F_{jump} + F_{sing}
\]
\end{example}

（绝对连续部分 + 跳跃部分 + 奇异部分，后两者的导数 a.e. 为零；$F_{sing}$ 的典型例子是 Cantor 函数型）。

\begin{example}[4（变差函数的连续性）]
$V(x) = V_a^x(F)$ 与 $F$ 同时左（右）连续：$F$ 在 $x_0$ 连续 $\iff$ $V$ 在 $x_0$ 连续。
\end{example}

\begin{example}[5（极大函数一般不可积）]
$Mf$ 仅弱 (1,1)，一般 $Mf \notin L^1$：取 $f(x) = \frac{1}{x(\ln x)^2}$（$0 < x < 1/2$ 型），$f \in L^1$ 但 $Mf(x) \gtrsim \frac{c}{x\vert\ln x\vert}$ 不可积。
\end{example}

\begin{example}[6（复合的绝对连续性）]
$F \in AC$、$\varphi \in C^1$，则 $\varphi \circ F \in AC$ 且链式法则 $(\varphi \circ F)' = \varphi' \circ F \cdot F'$ a.e. 成立。
\end{example}

\begin{quote}
\textbf{本章小结（核心图景）}：$\{\text{递增}\} \overset{Lebesgue}{\longrightarrow} \text{a.e. 可微}$；$BV \overset{Jordan}{=} \text{递增之差}$；$AC \subsetneq BV$；$\boxed{F = \int F' \iff F \in AC}$；微分定理把"平均"还原为"逐点"（Lebesgue 点、极大函数、Vitali 引理）；积分技巧（换元/分部/R–S）皆由 N–L 派生。
\end{quote}

\chap{第六章　Lp 空间}

收官章：把前五章的可积函数组织成赋范空间 $L^p$。主线：范数与基本不等式（Hölder、Minkowski）$\to$ 完备性（Riesz–Fischer）$\to$ 稠密性（简单函数、$C_c$、磨光逼近）$\to$ $L^2$ 的 Hilbert 空间结构与投影定理 $\to$ 各种收敛之间的关系 $\to$ 对偶空间 $(L^p)^\ast$ $\to$ 卷积与 Young 不等式，最后给出全课程总结。

\secn{6.1 定义与范数}

\begin{definition}[$L^p$ 范数]
设 $E$ 可测，$f$ 可测：
\end{definition}

\begin{itemize}

\item $1 \le p < \infty$：$\Vert f \Vert_p = \left( \int_E \vert f \vert^p \, dm \right)^{1/p}$；

\item $p = \infty$：$\Vert f \Vert_\infty = \inf\{ M : \vert f \vert \le M \text{ a.e.}\}$（本性上确界）。

\end{itemize}

$L^p(E) = \{ f : \Vert f \Vert_p < \infty \}$（对 a.e. 相等的函数取等价类，使 $\Vert f\Vert_p = 0 \iff f = 0$）。

\begin{proposition}[Young 不等式]
$a, b \ge 0$，$1 < p, q < \infty$，$\frac{1}{p} + \frac{1}{q} = 1$，则

\[
ab \le \frac{a^p}{p} + \frac{b^q}{q}
\]
\end{proposition}

（由 $\ln$ 的凹性/凸性，等号 $\iff$ $a^p = b^q$）。

\begin{theorem}[6.1（Hölder 不等式）]
$f \in L^p$，$g \in L^q$（$p, q$ 共轭），则 $fg \in L^1$ 且

\[
\int_E \vert f g\vert \, dm \le \Vert f \Vert_p \, \Vert g \Vert_q .
\]
\end{theorem}

\begin{proof}
不妨 $\Vert f\Vert_p, \Vert g\Vert_q > 0$；齐次化后设 $\Vert f \Vert_p = \Vert g \Vert_q = 1$，对 $\vert f\vert \vert g\vert$ 用 Young 不等式逐点相乘后积分。
\end{proof}

\begin{theorem}[6.2（Minkowski 不等式）]
$f, g \in L^p$（$1 \le p \le \infty$），则

\[
\Vert f + g \Vert_p \le \Vert f \Vert_p + \Vert g \Vert_p .
\]
\end{theorem}

\begin{proof}
$p = \infty$ 显然；$1 \le p < \infty$ 时由 $\vert f + g\vert^p \le 2^{p-1}(\vert f\vert^p + \vert g\vert^p)$ 知 $f + g \in L^p$，再用 $\vert f + g\vert^{p-1} \in L^{q}$ 与 Hölder 展开 $\int \vert f+g\vert^p$。

于是 $\Vert\cdot\Vert_p$ 是范数，$L^p$ 是赋范线性空间。
\end{proof}

\secn{6.2 完备性}

\begin{theorem}[6.3（Riesz–Fischer）]
$L^p(E)$（$1 \le p \le \infty$）是完备的，即 Banach 空间。
\end{theorem}

\begin{proof}
设 $\{f_n\}$ 是 Cauchy 列。取子列 $\{f_{n_k}\}$ 使

\[
\Vert f_{n_{k+1}} - f_{n_k} \Vert_p < 2^{-k},
\]

令 $g_m = \sum_{k=1}^{m} \vert f_{n_{k+1}} - f_{n_k}\vert$，$g = \lim_m g_m$。由 MCT 与三角不等式 $\Vert g_m\Vert_p \le \sum_{k \le m} 2^{-k} \le 1$，再由 Fatou 得 $g \in L^p$（从而 g 有限 a.e.）——级数 $\sum_k (f_{n_{k+1}} - f_{n_k})$ 绝对收敛 a.e.，定义

\[
f = f_{n_1} + \sum_{k=1}^{\infty} (f_{n_{k+1}} - f_{n_k}) \quad (\text{a.e. 收敛}).
\]

则 $f_{n_k} \to f$ a.e. 且 $\vert f\vert \le \vert f_{n_1}\vert + g \in L^p$；对 $\vert f_{n_k} - f\vert^p$ 用 Fatou：

\[
\Vert f_{n_k} - f \Vert_p^p = \int \lim_j \vert f_{n_k} - f_{n_j}\vert^p \le \liminf_j \Vert f_{n_k} - f_{n_j} \Vert_p^p \to 0.
\]

最后由 Cauchy 性把子列收敛放大到全列：$\Vert f_n - f\Vert_p \le \Vert f_n - f_{n_k}\Vert_p + \Vert f_{n_k} - f\Vert_p \to 0$。

（$p = \infty$：Cauchy 列在零测集外一致收敛，极限 a.e. 有界，收敛自动是 $L^\infty$ 的。）
\end{proof}

\begin{example}[有限测度时的包含关系]
$m(E) < \infty$、$q > p$ 时 $L^q \subset L^p$（Hölder：$\Vert f\Vert_p \le \Vert f\Vert_q\, m(E)^{\frac1p - \frac1q}$）；一般集合上各 $L^p$ 互不包含（如 $f(x) = x^{-1/p}$ 于 $(0,1)$ 与 $(1,\infty)$ 的对比）。
\end{example}

\secn{6.3 稠密性与逼近}

\begin{theorem}[6.4]
简单函数在 $L^p(E)$ 中稠密（$1 \le p < \infty$）。
\end{theorem}

\begin{proof}
$f = f^+ - f^-$，分别取递增简单列 $\varphi_n \nearrow f^\pm$（第三章），则 $\vert \varphi_n - f\vert^p \to 0$ 逐点且被 $\vert f\vert^p$ 控制，由 DCT 得 $\Vert \varphi_n - f\Vert_p \to 0$；有限测度集上截断处理无界性。
\end{proof}

\begin{theorem}[6.5]
$C_c(\mathbb{R}^n)$（紧支撑连续函数）在 $L^p(\mathbb{R}^n)$ 中稠密（$1 \le p < \infty$）。
\end{theorem}

\begin{proof}
先截断到有限测度集；简单函数用 Lusin 定理（闭集上连续 + 小测度余集）与正则性逼近；连续函数再用 Urysohn 型延拓/平移磨光。
\end{proof}

\begin{example}[磨光逼近]
标准磨光子 $\rho_\varepsilon(x) = \varepsilon^{-n}\rho(x/\varepsilon)$（$\rho \in C_c^\infty$，$\rho \ge 0$，$\int \rho = 1$），则

\[
\Vert \rho_\varepsilon \ast f - f \Vert_p \to 0 \quad (1 \le p < \infty),
\]

故 $C^\infty \cap L^p$ 也在 $L^p$ 中稠密（证明见 6.7 卷积的 Young 不等式 + 连续函数逼近）。
\end{example}

\begin{remark}
$p = \infty$ 时简单函数不稠密（一致极限保持"取值离散"的结构的闭包只能到简单函数类）。
\end{remark}

\secn{6.4 L² 空间与内积}

\begin{definition}[内积]
$\langle f, g \rangle = \int_E f\,\overline{g}\,dm$；诱导范数 $\Vert f \Vert_2^2 = \langle f, f \rangle$。
\end{definition}

\begin{proposition}[Cauchy–Schwarz 不等式]
$\vert\langle f, g\rangle\vert \le \Vert f \Vert_2\, \Vert g \Vert_2$（Hölder 在 $p = q = 2$ 的特例）。
\end{proposition}

由 Riesz–Fischer，$L^2$ 是完备内积空间，即 \textbf{Hilbert 空间}。

\begin{definition}[正交系]
$\{\varphi_k\} \subset L^2$：$\langle \varphi_j, \varphi_k \rangle = 0$（$j \ne k$）；若还有 $\Vert \varphi_k\Vert_2 = 1$，称为标准正交系。
\end{definition}

\begin{theorem}[6.6（投影定理）]
$M$ 是 $L^2$ 的闭子空间，$f \in L^2$，则存在唯一的 $g \in M$ 使

\[
\Vert f - g \Vert_2 = \inf_{h \in M} \Vert f - h \Vert_2,
\]
\end{theorem}

且 $f - g \perp M$（正交分解 $f = g + (f - g)$）。

\begin{proof}
存在性：取极小化序列 $h_n$，由平行四边形律

\[
\Vert h_n - h_m \Vert^2 = 2\Vert f - h_n\Vert^2 + 2\Vert f - h_m\Vert^2 - 4\left\Vert f - \tfrac{h_n + h_m}{2}\right\Vert^2 \to 0
\]

（$\frac{h_n+h_m}{2} \in M$，下确界性质使末项不小于 $4d^2$），故 $\{h_n\}$ Cauchy，极限 $g \in M$（闭子空间）取到下确界。正交性与唯一性由勾股定理/内积的严格正定性。
\end{proof}

\begin{remark}[Fourier 级数的预备]
$\{e^{inx}\}_{n \in \mathbb{Z}}$ 是 $L^2(-\pi, \pi)$ 的标准正交基（完备系），Fourier 级数 = 投影定理的展开——详见调和分析课程。
\end{remark}

\secn{6.5 各种收敛的关系}

设 $\{f_n\}$、$f$ 是 a.e. 有限的可测函数：

\begin{itemize}

\item \textbf{一致收敛 $\Rightarrow$ 处处收敛 $\Rightarrow$ a.e. 收敛}；一致收敛 + 有限测度 $\Rightarrow$ 依测度收敛；

\item \textbf{a.e. 收敛 $+ \, m(E) < \infty$ $\Rightarrow$ 依测度收敛}（Egorov/Lebesgue 定理，第三章）；

\item \textbf{依测度收敛 $\not\Rightarrow$ a.e. 收敛}，但 \textbf{Riesz}：可抽子列 a.e. 收敛；

\item \textbf{$L^p$ 收敛 $\Rightarrow$ 依测度收敛}（Chebyshev：$m\{\vert f_n - f\vert \ge \sigma\} \le \sigma^{-p}\Vert f_n - f\Vert_p^p$）$\Rightarrow$ 可抽子列 a.e. 收敛；

\item \textbf{依测度 / $L^p$ 收敛都不推出 a.e. 收敛}（滑脊例）；

\item $L^p$ 收敛 $\Rightarrow$ 范数收敛 $\Vert f_n\Vert_p \to \Vert f\Vert_p$（三角不等式）。

\end{itemize}

\begin{example}[滑脊函数]
$f_n = n^{1/p}\,\chi_{(0,\, 1/n)}$：$\Vert f_n \Vert_p = 1$，$f_n \to 0$ 依测度（也在 $(0,1]$ 上逐点趋于 0），但 $\Vert f_n - 0\Vert_p = 1 \not\to 0$——依测度收敛不蕴含 $L^p$ 收敛。
\end{example}

\begin{example}[反向滑脊]
$f_n = \chi_{(\frac{1}{n+1},\, \frac{1}{n}]}$ 型：$\int \vert f_n\vert^p \to 0$（$L^p \to 0$），$f_n \to 0$ 依测度，但对每个 $x \in (0,1]$，$f_n(x) = 1$ 无穷多次——不 a.e. 收敛（只有子列收敛）。
\end{example}

\secn{6.6 对偶空间}

\begin{theorem}[6.7]
设 $1 < p < \infty$，$q$ 为共轭指标，则 $(L^p)^\ast \cong L^q$：每个连续线性泛函 $L \in (L^p)^\ast$ 都唯一对应 $g \in L^q$ 使

\[
L(f) = \int_E f\,g\,dm, \qquad \Vert L \Vert = \Vert g \Vert_q .
\]
\end{theorem}

\begin{proof}[证明思路]
"$\le$" 是 Hölder。构造 $g$（$1 < p < \infty$，$\sigma$-有限情形）：

Step 1：先设 $m(E) < \infty$。定义 $\nu(A) = L(\chi_A)$（$A$ 可测），则 $\nu$ 是有限符号测度且 $\nu \ll m$（$m(A) = 0 \Rightarrow \chi_A = 0$ in $L^p$ $\Rightarrow \nu(A) = 0$）；

Step 2：Radon–Nikodym 定理给出 $g = \frac{d\nu}{dm}$；

Step 3：$L(f) = \int f g$ 对简单函数成立，由稠密性（6.3）与 $\vert L(f)\vert \le \Vert L\Vert \Vert f\Vert_p$ 推广到全体 $L^p$，并得 $\Vert g\Vert_q \le \Vert L\Vert$（取 $f = \vert g\vert^{q-2}\overline{g}$ 型检验函数）。
\end{proof}

\begin{remark}
（i）$(L^1)^\ast \cong L^\infty$（同样成立）；（ii）$(L^\infty)^\ast \supsetneq L^1$——除积分泛函外还有"在零测集外看不见"的奇异泛函（Banach 极限型），其刻画需泛函分析工具（Hahn–Banach）；（iii）Radon–Nikodym 定理的证明用 Hahn 分解，见一般测度论教材。
\end{remark}

\secn{6.7 卷积与 Young 不等式}

\begin{definition}[卷积]
$f \in L^p$，$g \in L^1$：

\[
(f \ast g)(x) = \int_{\mathbb{R}^n} f(x - y)\, g(y)\, dy .
\]
\end{definition}

（被积函数对 a.e. 的 $x$ 绝对可积，$f \ast g$ 良定义且 a.e. 有限。）

\begin{theorem}[6.8（Young 不等式）]
$\Vert f \ast g \Vert_p \le \Vert f \Vert_p \, \Vert g \Vert_1$（$1 \le p \le \infty$）。
\end{theorem}

\begin{proof}
$f \ast g(x) = \int f(x-y)g(y)\,dy$ 视为"关于 $y$ 的积分的三角不等式"（Minkowski 积分形式）：

\[
\Vert f \ast g \Vert_p = \left\Vert \int f(\cdot - y)\,g(y)\,dy \right\Vert_p \le \int \Vert f(\cdot - y)\Vert_p \,\vert g(y)\vert\,dy = \Vert f\Vert_p \Vert g\Vert_1. \; \]
\end{proof}

\begin{proposition}[卷积的初等性质]
交换律 $f \ast g = g \ast f$、结合律、与平移交换 $\tau_h(f\ast g) = (\tau_h f)\ast g = f \ast (\tau_h g)$；$f \in L^p$、$g \in C_c$ 时 $f \ast g$ 连续（$p < \infty$），$g \in C_c^\infty$ 时 $f \ast g \in C^\infty$ 且求导可移到 $g$ 上。
\end{proposition}

\begin{example}[磨光逼近的证明]
$\rho_\varepsilon \ast f - f = \int [f(\cdot - y) - f]\,\rho_\varepsilon(y)\,dy$，对 $C_c$ 中的 $f$ 由一致连续性控制，一般 $f \in L^p$ 用 6.3 的稠密性 + Young 不等式：

\[
\Vert \rho_\varepsilon \ast f - f \Vert_p \le \Vert \rho_\varepsilon \ast (f - \varphi) \Vert_p + \Vert \rho_\varepsilon \ast \varphi - \varphi \Vert_p + \Vert \varphi - f \Vert_p \to 0. \; \]
\end{example}

\textbf{应用.} 偏微分方程的正则化：用磨光子逼近非光滑初值/系数，得到光滑化近似解。

\secn{6.8 全课程总结}

\textbf{六章主线}：

第一章 集合论基础（动机：Riemann 积分三大缺陷）$\to$ 第二章 Lebesgue 测度（外测度、Carathéodory 可测、结构定理）$\to$ 第三章 可测函数（运算、逼近、Egorov/Lusin、依测度收敛）$\to$ 第四章 Lebesgue 积分（MCT/Fatou/DCT、Fubini）$\to$ 第五章 微分与 N–L 公式（BV/AC 链、Lebesgue 微分定理）$\to$ 第六章 $L^p$ 空间（完备性、稠密性、对偶、卷积）。

\textbf{核心定理清单}：

\\begin{center}
\\begin{tabular}{cl}
\\toprule
章 & 核心定理 \\
\\midrule
1 & Cantor--Bernstein；$\mathbb{R}$ 不可数（对角线法） \\
2 & Carath\'eodory 扩张定理；结构定理；Vitali 不可测集 \\
3 & 简单函数逼近；Egorov；Lusin；Riesz 子列定理 \\
4 & Levi 单调收敛；Fatou；控制收敛；Tonelli/Fubini \\
5 & Lebesgue 微分定理；Jordan 分解；N--L 公式 $\iff$ AC \\
6 & H\"older/Minkowski；Riesz--Fischer；$(L^p)^\ast = L^q$；Young \\
\\bottomrule
\\end{tabular}
\\end{center}

\textbf{后续课程衔接}：泛函分析（Banach/Hilbert 空间，$L^p$ 是基本例子）、调和分析（卷积、Fourier 变换、极大函数）、偏微分方程（Sobolev 空间 $W^{k,p}$）、概率论（期望与极限定理的 $L^p$ 语言）。

\secn{6.9 习题课}

\begin{exercise}[1]
用 Young 不等式证明 Hölder 不等式，并给出等号条件：等号 $\iff$ $\frac{\vert f\vert^p}{\Vert f\Vert_p^p} = \frac{\vert g\vert^q}{\Vert g\Vert_q^q}$ a.e.
\end{exercise}

\begin{exercise}[2]
（内插）设 $m(E) < \infty$，$f \in L^{p_1} \cap L^{p_2}$（$p_1 < p_2$），证明：对一切 $p \in (p_1, p_2)$ 有 $f \in L^p$，且

\[
\Vert f \Vert_p \le \Vert f \Vert_{p_1}^{\theta}\, \Vert f \Vert_{p_2}^{1-\theta}, \qquad \frac{1}{p} = \frac{\theta}{p_1} + \frac{1-\theta}{p_2}.
\]
\end{exercise}

（Hölder 应用于 $\vert f\vert^{\theta p}$ 与 $\vert f\vert^{(1-\theta)p}$——Riesz–Thorin 内插的雏形。）

\begin{exercise}[3]
证明：$f_n \to f$ 于 $L^p$ $\Rightarrow$ $\Vert f_n\Vert_p \to \Vert f\Vert_p$ 且 $f_n \Rightarrow f$（依测度）；构造反例说明两个必要条件合起来仍不充分。
\end{exercise}

\begin{exercise}[4]
（稠密性练习）验证阶梯函数、简单函数、$C_c$ 在 $L^p(\mathbb{R}^n)$（$p < \infty$）中层层稠密。
\end{exercise}

\begin{exercise}[5]
（卷积练习）证明 $\Vert \rho_\varepsilon \ast f - f\Vert_p \to 0$（$f \in L^p$，$1 \le p < \infty$）。
\end{exercise}

\begin{exercise}[6]
（对偶计算）求泛函 $L(f) = \int f(x)\,h(x)\,dx$ 在 $L^2$ 上的范数（$h$ 给定）：$\Vert L\Vert = \Vert h\Vert_2$。
\end{exercise}

\begin{remark}[考试重点]
Hölder/Minkowski 的证明、Riesz–Fischer 完备性证明、简单函数/$C_c$ 稠密性证明、N–L 公式的证明。
\end{remark}

\begin{quote}
\textbf{全书完}：六章手写笔记电子化完成。记号与证明主线以课堂笔记为准，个别习题与反例细节参考了标准教材（Royden、周民强《实变函数论》）。
\end{quote}


\end{document}
